\subsection{Modular Invariant Starobinsky Inflation and the Species Scale --- arXiv:2407.12081}

\textbf{Authors:} Gonzalo F. Casas, Luis E. Ibáñez

\paragraph{主な主張}
species scale $\Lambda$ の勾配から $V\sim G^{i\bar j}(\partial_i\Lambda)(\partial_{\bar j}\Lambda)/\Lambda^2$ を inflaton potential として構成し、単一の $SL(2,\mathbb{Z})$ modulus では large-modulus 極限で Starobinsky 型の plateau と $N_e\sim\mathcal N$, $\epsilon\sim\Lambda^4$ を得ることを示した \cite{Casas:2024species}。

\paragraph{新規性}
具体的な superpotential からではなく、duality invariant な species scale とその微分を potential の構成原理として用いることで、tower of states・量子重力 cutoff と modular inflation の観測量を直接結び付けた点が新しい。

\paragraph{位置付け}
microscopic な string/SUGRA 導出は未完成で D-term 起源は可能性として提示されるに留まるが、species scale を基礎量として modular invariant inflation を構成する代表的な bottom-up proposal である。
